\input{\string~/Documents/School/"UC Irvine"/ta_template2.0.tex}
\rh{2B}{7.1}
\chead{\today}
\graphicspath{ {\string~/Desktop} }
\usetikzlibrary{arrows}
%============ANSWER KEY TOGGLE===========
\newcommand{\ans}[1]{ \noindent {\color{red} \textbf{Answer:} #1}}
%\renewcommand{\ans}[1]{\vspace{3in}}
%==========================================
\begin{document}
\begin{center}
\textbf{Integration by Parts}
\end{center}

\begin{prop}[Integration by Parts]
We have the following integration identity:
\[
\int u dv = uv - \int v du
\]
\end{prop}
\begin{aun}
Here's a cool rule of thumb that I found to always select the correct $u$ and $v$: \textbf{Choose $u$ in this order:}
\begin{description}
\item[L:] logarithms
\item[I:] inverse functions
\item[A:] algebraic (polynomials, etc.)
\item[T:] trigonometric functions
\item[E:] exponentials
\end{description}
\end{aun}

\begin{aun}
I found out later that if we add a $d$ at the end, this spells DETAIL backwards... so another way to remember this trick is:
\begin{description}
\item[D:] $dv$
\item[E:] exponentials
\item[T:] trigonometric functions
\item[A:] algebraic (polynomials, etc.)
\item[I:] inverse functions
\item[L:] logarithms
\end{description}
\end{aun}


\problembox{[Introductory Problems]
Compute the antiderivative:
\begin{enumerate}[label=\alph*.]
\item  $ \displaystyle \int (x^2 + 2x)\cos(x) dx$
\item $\displaystyle \int 4x \cos (2-3x) dx$
\end{enumerate}
}
\ans{
%TEACH THE TABULAR METHOD
\begin{enumerate}[label=\alph*.]
\item We show the following steps:
\begin{alignat*}{2}
	\int (x^2 + 2x)\cos(x) dx &= \int (x^2 + 2x) d \sin(x) \\
					&= \lrp{x^2 + 2x} \sin(x) - \int \sin(x) d x^2 + 2x\\
					&= \lrp{x^2 + 2x} \sin(x) - 2 \int \sin(x) (x+1) dx\\
					&= \lrp{x^2 + 2x} \sin(x) - 2 \int (x+1) d \lrp{-\cos(x)}\\
					&= \lrp{x^2 + 2x} \sin(x) + 2 \lrp{(x+1) \cos(x) - \int \cos(x)dx}\\
					&= \lrp{x^2 + 2x} \sin(x) + 2 \lrp{(x+1) \cos(x)} -\sin(x) + C
\end{alignat*}
\item Start with a change of coordinates:
\[
u = 2-3x \mte{ , } du = -3dx \mte{ and } x = \frac{u - 2}{-3}
\]
From here, we can rewrite:
\[
\int 4x \cos(2-3x) dx = \frac{-4}{3} \int \underbrace{x}_{= \frac{u - 2}{-3}} \cos (\underbrace{2 - 3x}_{u}) \underbrace{-3x dx}_{du} = \frac{-4}{3} \cdot \frac{1}{-3} \int (u-2) \cos (u) du
\]
which can be solved the same way as the preceding problem.
\end{enumerate}
}

\problembox{[Interesting Twists] Compute the antiderivative:
\begin{enumerate}
\item $\displaystyle \int e^x \sin (x) dx$
\item  $ \displaystyle\int \ln\sqrt{x} dx$
\item $\displaystyle \int e^{\sqrt{x}} dx$
\item $ \displaystyle \int \cos(\ln(x))dx$
\item $\displaystyle \int (\arcsin(x))^2 dx$
\end{enumerate}
}
\ans{
We can use tabular integration to speed these up.
\begin{enumerate}
\item Consider the following table:
\[
\begin{array}{|c|c|}
\hline
\frac{d}{dx} & \int dx \\ \hline
\sin(x) & e^x \\
\cos(x) & e^x\\
-\sin(x) & e^x\\
-\cos (x) & e^x \\ \hline
\end{array}
\]
This gives us:
\begin{alignat*}{2}
\int e^x \sin(x)dx &= e^x \sin(x) - \int e^x d(\sin (x)) \\
			&= e^x \sin(x) - \int e^x \cos(x)dx\\
			&= e^x \sin(x) - \lrb{e^x \cos(x) - \int e^x d(\cos(x))}\\
\underbrace{\int e^x \sin(x)dx } &= e^x \sin(x) - e^x \cos(x) - \underbrace{\int e^x \sin(x) dx}
\end{alignat*}
Looking at the last line, we see that:
\[
\int e^x \sin(x) dx = \frac12 e^x (\sin(x) - \cos (x))
\]
\item Using $u$-substitution, we see that:
\[
u = \sqrt{x} \mte{ and } du = \frac{1}{2 \sqrt{x}} dx \mte{ give us } dx = 2 u du
\]
This converts our integral to
\[
\int \ln(\sqrt{x}) = 2 \int u \ln (u) du
\]
which we can compute using the dETAIL technique.
\item Same technique as above gives us:
\[
\int e^{\sqrt{x}} dx = 2 \int ue^u du
\]
\item Making a $u$-substitution for $\ln(x)$ gives you:
\[
\underbrace{u = \ln(x)}_{e^u = x} \mte{ and } du = \frac1x dx \mte{ implying } dx = e^u du
\]
This turns the integral into:
\[
\int\cos (\ln(x))dx = \int e^u \cos(u) du
\]
which we computed in the example above. 
\item This is another $u$-substitution, this time:
\[
\underbrace{u = \arcsin (x)}_{x = \sin(u)} \mte{ and } du = \frac{1}{\sqrt{1-x^2}} dx
\]
Solving for $u$ in the second equation, we see:
\[
dx = \sqrt{1 - (\sin(u)^2)} du = \cos(u) du
\]
by the trigonometric identity $\sin^2(x) + \cos^2(x) = 1$. This gives us the integral:
\[
\int \arcsin^2 (x) dx = \int u^2 \cos(u) du
\]
for which a similar integral was computed earlier. 
\end{enumerate}
}

\end{document} 